\chapter{Unplanned pregnancy}
\index{Unplanned pregnancy}

\section*{Definition and Overview}
An unplanned pregnancy is a pregnancy that occurs when it was not intended, either because contraception was not used, failed, or was used incorrectly, or because conception happened before a decision to become pregnant. It is distinct from an unwanted pregnancy, although the terms overlap. Some unplanned pregnancies are accepted and welcomed once they occur, while others are deeply unwanted. Around half of all pregnancies worldwide are estimated to be unplanned, making this one of the most common reproductive health experiences a person can face. An unplanned pregnancy raises important questions about options, timing, relationships, finances, and health, and people deserve non-judgemental information and support to make the decision that is right for them. In many countries, the three main options for an unplanned pregnancy are continuing the pregnancy and raising the child, continuing the pregnancy and arranging adoption or foster care, and ending the pregnancy through abortion, which may be medical or surgical. Access to these options varies by local, so accurate, local, up-to-date advice is essential.

\section*{Epidemiology}
Unplanned pregnancy is common worldwide. Globally, estimates suggest that roughly 44\% of all pregnancies are unintended. Rates vary widely between countries and between population groups within a country. Younger people, those with lower incomes or education, people in unstable relationships, and those using less effective contraception are at higher risk. In many countries, although overall rates have declined in recent decades alongside improved access to effective long-acting contraception, a substantial proportion of pregnancies each year are still reported as unintended. Teenage pregnancy rates in many countries have fallen considerably, but older women experiencing contraceptive failure or changes in fertility also account for many unplanned pregnancies. Condom failure and missed or forgotten contraceptive pills are among the most common reasons cited. Emergency contraception is widely used, but many people either do not know about it, cannot access it in time, or do not realise it is available without a prescription.

\section*{Aetiology and Pathophysiology}
An unplanned pregnancy occurs when sperm fertilises an egg and the resulting embryo implants in the uterus, despite the pregnancy not being intended. The immediate causes relate to contraception and timing.

\begin{itemize}
    \item \textbf{No contraception:} not using any method, often because of cost, access, side effects, misinformation, or ambivalence
    \item \textbf{Contraceptive failure:} even perfect use is not 100\% effective, for example with typical-use failure rates of the combined pill, condoms, and the injection
    \item \textbf{Incorrect or inconsistent use:} missed pills, expired or broken condoms, incorrect diaphragm fitting, or late depot injections
    \item \textbf{Emergency contraception failure:} using an oral emergency contraceptive too late or after further unprotected sex
    \item \textbf{Sexual coercion or assault:} pregnancy resulting from non-consensual sex
    \item \textbf{Withdrawal or fertility awareness errors:} relying on methods with high failure rates or on incorrectly calculated fertile days
    \item \textbf{Breastfeeding assumptions:} believing that breastfeeding fully prevents ovulation when it does not reliably do so
\end{itemize}

Underlying physiological factors matter too. Ovulation and fertility can return earlier than expected after childbirth or stopping a contraceptive. Irregular cycles make the fertile window hard to predict. Tubal ligation and vasectomy are highly effective but can rarely fail. Contraceptive implants and intrauterine devices (IUDs) are the most effective reversible methods, but even these have a small failure rate, and an unrecognised expulsion of an IUD allows fertility to return.

\section*{Clinical Features}
The presenting feature of an unplanned pregnancy is usually a missed period, but the diagnosis may come sooner or later depending on circumstances.

\begin{itemize}
    \item A missed or delayed period in a person who usually menstruates regularly
    \item A positive home pregnancy test taken for any reason
    \item Breast tenderness, nausea, fatigue, or frequent urination
    \item Light spotting around the time the period was expected
    \item The pregnancy being discovered only at a routine appointment, such as a cervical screening test or a contraception review
    \item Emotional responses that range widely, including shock, denial, distress, guilt, or relief, and these feelings may change over time
\end{itemize}

A pregnancy test is the cornerstone of diagnosis. Modern home urine tests are highly accurate from the first day of a missed period. A person who has irregular periods, is using hormonal contraception, is breastfeeding, or is around menopause may not notice a missed period, so the pregnancy may be detected later.

\section*{Differential Diagnosis}
\begin{itemize}
    \item A false-negative home pregnancy test taken too early or from dilute urine
    \item A false-positive home pregnancy test, which can occasionally occur with certain medications, recent pregnancy loss, or an incorrectly read test
    \item A non-viable pregnancy, such as an early miscarriage or blighted ovum, which will show a positive test initially
    \item An ectopic pregnancy, which also produces a positive test but requires urgent assessment
    \item A missed or delayed period from other causes, including stress, weight change, polycystic ovary syndrome, thyroid disease, or perimenopause
    \item Pituitary or other hormone problems that produce pregnancy-like symptoms
\end{itemize}

Whenever a home test is positive, the diagnosis should be confirmed in a clinical setting, especially when the person needs to make decisions about options or when symptoms are unusual. A negative home test with persistent symptoms should also be checked, because tests taken very early can be wrong.

\section*{When to Seek Medical Care}
Anyone with a positive pregnancy test should seek care promptly, whether they plan to continue the pregnancy or not, so that they can confirm the pregnancy, date it, discuss options, and arrange appropriate care. People should also see a doctor if a pregnancy test is unclear or if they have been unable to access contraception and are worried they may be pregnant.

\textbf{Contact your local emergency services for:}
\begin{itemize}
    \item Severe one-sided lower abdominal pain together with a positive pregnancy test, which may indicate an ectopic pregnancy
    \item Heavy vaginal bleeding with clots and pain, which may indicate a miscarriage or ectopic pregnancy
    \item Fainting, dizziness, or collapse with shoulder-tip pain
    \item Fever with pelvic pain in early pregnancy
    \item Any symptoms that feel life-threatening or unbearable
\end{itemize}

\section*{Diagnosis and Investigations}
The assessment of a person with a suspected unplanned pregnancy is designed to confirm the pregnancy, estimate how advanced it is, and identify any immediate health concerns.

\begin{itemize}
    \item \textbf{Pregnancy test:} a urine or blood test for human chorionic gonadotrophin (hCG) confirms the pregnancy
    \item \textbf{History:} the date of the last menstrual period, cycle regularity, contraceptive use, symptoms, past pregnancies, medical conditions, and any bleeding or pain
    \item \textbf{Physical examination:} blood pressure, heart rate, and gentle abdominal examination; a pelvic examination is not always necessary in early pregnancy
    \item \textbf{Ultrasound:} performed from about five to six weeks of pregnancy to confirm a heartbeat, date the pregnancy, and locate it in the uterus, which is essential to exclude an ectopic pregnancy
    \item \textbf{Blood tests:} blood group, Rh status, full blood count, and screening for sexually transmitted infections where relevant
    \item \textbf{Blood hCG measurements over time:} useful when the pregnancy is very early or when ultrasound is inconclusive
\end{itemize}

The dating of the pregnancy is critical because the options available, and the timing of those options, depend on gestational age. In many countries, surgical and medical abortion are available under different regulations in different states and territories, so local advice on access and timeframes should always be sought.

\section*{Management}
Management of an unplanned pregnancy is about supporting the person's decision, whatever it is. There is no right answer for everyone, and the health system's role is to provide accurate information, compassionate care, and safe options.

\begin{itemize}
    \item \textbf{Continuing the pregnancy:} referral for antenatal care, including blood tests, ultrasound dating, folic acid, and advice on nutrition, alcohol, smoking, and medications
    \item \textbf{Parenting support:} access to parenting programmes, financial support, housing, and community services for those who choose to raise the child
    \item \textbf{Adoption and foster care:} counselling and support through authorised agencies for those considering alternative care arrangements
    \item \textbf{Medical abortion:} medicines (usually mifepristone followed by misoprostol) used in early pregnancy, either at home with support or in a clinic, depending on gestational age and local availability
    \item \textbf{Surgical abortion:} a suction or surgical procedure under local or general anaesthetic, offered within time limits that vary by local
    \item \textbf{Psychological support:} counselling, whether before, during, or after a decision, because the emotional impact can be significant and long-lasting
    \item \textbf{Contraception:} arranging effective contraception to start immediately after any procedure or at the appropriate time, so that future pregnancies are planned if desired
\end{itemize}

Whatever the decision, it should be made freely, without pressure from partners, family, or clinicians, and within the limits of the law in the person's local. Clinicians should provide unbiased information and support referrals to services that can assist.

\section*{Complications}
\begin{itemize}
    \item Unrecognised or late presentation of pregnancy, delaying antenatal care and increasing pregnancy-related risks
    \item Ectopic pregnancy, which is a medical emergency and requires treatment
    \item Unsafe abortion attempts when safe services are inaccessible, which carry risks of severe bleeding, infection, and death
    \item Psychological distress, including anxiety, depression, and relationship conflict, whatever the decision
    \item Financial and social strain for those continuing an unplanned pregnancy without support
    \item Complications of the chosen option, such as heavy bleeding or infection after an abortion, which are uncommon when care is provided safely
\end{itemize}

\section*{Prognosis and Outcomes}
With safe, timely care, the physical outcomes of an unplanned pregnancy, whatever the decision, are generally good. Continuing the pregnancy with early antenatal care has outcomes similar to those of planned pregnancies for most people. Medical and surgical abortion in many countries are very safe when performed by trained providers within legal limits, with serious complications being rare. The emotional and social outcomes depend heavily on the support available and whether the decision felt free and informed. People who feel supported in their choice, and who receive contraception and follow-up afterwards, tend to have better outcomes. Reassuringly, most people who have an abortion do not regret it in the long term, although it is normal to have mixed feelings.

\section*{Prevention and Self-Care}
\begin{itemize}
    \item Choose an effective contraceptive and use it consistently; long-acting methods such as the implant and IUD are the most reliable
    \item Know that typical-use failure rates of the pill and condoms are higher than most people think
    \item Use condoms as well as another method to prevent pregnancy and sexually transmitted infections
    \item Know how to access emergency contraception and use it as soon as possible after unprotected sex; it works best the sooner it is taken
    \item After an abortion or childbirth, ovulation can return before the first period, so start contraception immediately if you do not want another pregnancy
    \item Seek prompt medical review for any positive pregnancy test so that you can make an informed decision in good time
    \item Reach out for counselling or support services; you do not have to make this decision alone
    \item Be kind to yourself; feelings of shock, guilt, or sadness are normal, and professional support is available
\end{itemize}

\section*{Sources and Further Reading}
The following organisations provide reliable information on unplanned pregnancy, options, and contraception.

\begin{itemize}
    \item Healthdirect Australia. \textit{Unplanned pregnancy and pregnancy options.} https://www.healthdirect.gov.au/
    \item MSD Manual. \textit{Overview of contraception and pregnancy options.} https://www.msdmanuals.com/
    \item MedlinePlus. \textit{Unintended pregnancy.} https://medlineplus.gov/
    \item NHS. \textit{Unplanned pregnancy: your choices.} https://www.nhs.uk/
    \item World Health Organization. \textit{Abortion care guideline.} https://www.who.int/
    \item Marie Stopes Australia. \textit{Pregnancy options counselling and services.} https://www.msaustralia.org.au/
\end{itemize}
